\section{Discussion}
\label{chap:discussion}

This chapter synthesizes the experimental results from Chapter~\ref{chap:experiments}, critically evaluating the performance and efficiency profiles of the alternative algorithms against their respective BP baselines. The analysis focuses on interpreting the observed trade-offs, positioning the findings within the existing literature, and assessing the practical implications of these alternative training paradigms, with a particular focus on energy efficiency. The evolutionary narrative of BP-free methods, progressing from the FF's foundational demonstration through CaFo's structured design to the emergence of MF as a highly efficient and performant alternative, provides the guiding framework for this discussion.

\subsection{Interpretation of Key Findings}
\label{sec:discussion_interpretation}

The empirical results presented in Chapter~\ref{chap:experiments} provide a nuanced perspective on the capabilities of BP-free algorithms, tracing an evolutionary path toward the dual objectives of competitive accuracy and superior efficiency.

Geoffrey Hinton's FF algorithm established a crucial proof of concept by validating that deep networks can attain competitive accuracy without a backward pass (Table~\ref{tab:ff_bp_mlp_perf_results_ch4}). However, this validation was performed at a prohibitive efficiency cost. As illustrated in Figure~\ref{fig:ch4_ff_bp_mnist_mlp_4x2000_conv_time}, FF's convergence process is markedly slower and more volatile than that of backpropagation, culminating in significantly extended training durations and elevated energy consumption. The analysis at the hardware level (Figure~\ref{fig:ch4_ff_bp_mnist_mlp_4x2000_resource_util}) critically revealed that FF's computational patterns result in suboptimal GPU utilization and that its presumed memory advantages did not materialize for the MLP architectures tested. This positions FF as a \textbf{conceptually significant yet practically unviable baseline}, thereby motivating the investigation of more sophisticated methods.

The CaFo algorithm signifies a structured advancement, designed to deliver a more direct supervised signal via its block-wise predictor architecture. The experimental findings revealed a clear trade-off between its two evaluated variants. CaFo-Rand-CE, which uses fixed random blocks, offered consistent, though modest, peak memory savings, but incurred a substantial accuracy penalty on complex datasets (Table~\ref{tab:cafo_bp_cnn_perf_results_ch4_reordered}). This result highlights that effective learning cannot be achieved simply by adding local classifiers to random feature projections. In contrast, CaFo-DFA-CE substantially improved accuracy by pre-training its blocks with DFA, approaching BP's performance. However, this success was accompanied by a considerable computational burden, as the DFA pre-training phase led to profoundly higher energy and time consumption (Table~\ref{tab:cafo_bp_cnn_eff_results_ch4_reordered_eff}). CaFo consequently presents a challenging decision: one must either accept a significant performance deficit for a niche efficiency gain or incur a heavy computational cost to achieve competitive accuracy.

The MF algorithm, the most advanced method investigated, emerges from these experiments as a particularly compelling alternative for its native MLP architectures. As presented in Table~\ref{tab:mf_bp_mlp_perf_ch4}, MF \textbf{consistently achieved or exceeded the accuracy of rigorously optimized BP baselines}. This enhanced generalization performance is attributed to its sequence of local optimizations converging to a more advantageous final validation loss compared to BP's global optimization strategy (Figure~\ref{fig:mf_bp_cifar10_mlp_conv_curves}). This superior performance coincided with profound efficiency improvements. On the CIFAR-10 MLP task, MF proved to be approximately 34\% faster and consumed approximately 41\% less energy than its BP counterpart (Table~\ref{tab:mf_bp_mlp_eff_ch4}). This efficiency is a direct consequence of its computationally lighter design, which eliminates the backward pass. MF thus effectively navigates the trade-off between accuracy and efficiency, \textbf{substantially fulfilling the promise of a practical, high-performance, and sustainable BP-free algorithm for MLPs}.

A critical methodological element supporting these interpretations is the consistent application of systematic hyperparameter Optuna optimization and validation-based early stopping. This protocol ensured that every algorithm was evaluated at its optimal performance point, facilitating a fair comparison of their intrinsic characteristics.

\subsection{Algorithm and Architecture Strengths and Weaknesses}
\label{sec:discussion_strengths_weaknesses}

Each algorithm demonstrated distinct strengths and weaknesses that were frequently related to its unique learning mechanism and its interaction with the native network architecture. The Forward-Forward (FF) algorithm's primary strength lies in its conceptual contribution: It fundamentally established that BP-free local learning can achieve competitive accuracy. However, its weaknesses are profound from a practical point of view, including a high inefficiency in both time and energy.

The Cascaded-Forward (CaFo) algorithm presents a more complex set of trade-offs. Its main weakness stems from a crucial principle revealed by these results: \textbf{the quality of feature representations is paramount in BP-free learning}. The mere attachment of local predictors, as in the Rand-CE variant, is insufficient for complex tasks. Conversely, while the DFA-CE variant produces effective features, its pre-training stage is exceptionally resource-intensive, negating any overall energy efficiency.

Mono-Forward (MF) on MLPs demonstrated the most compelling balance of strengths. It consistently achieved accuracy on par with or superior to fair BP baselines, coupled with significant reductions in training time and energy consumption. However, a critical and counterintuitive weakness emerged from the analysis: \textbf{The theoretical memory savings of BP-free learning did not fully materialize}. The peak memory advantage was modest or non-existent on the tested MLPs. This empirically demonstrates that the practical memory overhead of MF's unique components, specifically the learnable projection matrices and their optimizer states, \textbf{can partially or fully counteract the gains} from forgoing activation storage, challenging a common assumption about BP-free methods.

Finally, backpropagation (BP) serves as a benchmark for a highly mature, robust, and generally efficient algorithm, benefiting from decades of co-optimization. Its well-known weakness is the inherent memory and computational cost of the backward pass. This analysis highlights that the optimal approach is context-dependent, with MF exhibiting a particularly strong profile for applications based on MLPs.

\subsection{Practical Implications}
\label{sec:discussion_implications}

The results of this paper carry several practical implications for the development and deployment of deep learning models, particularly concerning sustainability. The most direct consequence is the potential for significant energy savings and a reduced environmental impact. By delivering substantial reductions in training time and energy consumption while enhancing accuracy, the MF algorithm \textbf{presents a concrete path toward more sustainable AI}, translating directly into a lower estimated CO2e. These efficiency gains also make such algorithms highly advantageous for deployment on edge devices or in other resource-constrained environments. The emergence of effective BP-free algorithms like MF \textbf{could also inspire new hardware architectures} optimized for their specific computational patterns, potentially altering trade-offs in future chip designs. Ultimately, this research provides empirical data to inform algorithm selection. For tasks based on MLPs, \textbf{MF appears to be a better choice than BP when both accuracy and efficiency are primary considerations}, suggesting that this direct local loss paradigm is a powerful direction for future algorithm development.

\pagebreak

\subsection{Theoretical Insights}
\label{sec:discussion_theoretical}

This comparative investigation also yields several theoretical insights into the mechanisms of deep learning. The results robustly demonstrate that \textbf{effective high-performance learning can be realized using purely local rules}, a foundational premise established by FF and compellingly validated by MF. The success of MF's mechanism suggests that complex credit assignment can be managed without global error propagation. By eschewing the precise symmetric weight transport of BP, these algorithms align more closely with the principles of biological learning. A key related insight is the strong coupling between the algorithm and its architecture, suggesting that \textbf{future progress in BP-free learning will likely depend on the co-design of both the learning rule and the network primitives} upon which it operates. Finally, MF's ability to converge to a lower final validation loss than BP on the CIFAR-10 MLP task (Figure~\ref{fig:mf_bp_cifar10_mlp_conv_curves}) offers a fascinating insight into optimization landscapes. It suggests that a sequence of greedy local optimizations may guide a model to a more advantageous region of the parameter space, \textbf{challenging the conventional wisdom that only global optimization can locate the best solutions}.

\subsection{Limitations and Critical Evaluation}
\label{sec:discussion_limitations}

Although this paper provides a rigorous comparison under controlled conditions, several limitations must be acknowledged. First, the author intentionally focused the experimental scope on the \textit{native} architectures for each algorithm. Consequently, the \textbf{strong performance of MF on MLPs does not automatically guarantee similar success on different architectures}, such as Transformers or deeper CNNs, without further research.

A second set of limitations arises from algorithm-specific nuances. The inefficiency of FF makes its current formulation impractical. For CaFo, the DFA-CE variant proved too resource-intensive to be "efficient", while the Rand-CE variant's accuracy remains a significant weakness. Most importantly, although highly successful on MLPs, \textbf{MF's peak memory savings were modest}, a critical nuance highlighting that the overhead of its unique components can counteract a significant portion of the gains from eliminating the backward pass.

Furthermore, while systematic hyperparameter optimization was conducted, such searches are not exhaustive. The findings are also tied to a specific hardware and software stack, and the estimated CO2e depends on the local grid's carbon intensity. Finally, this paper focuses primarily on training efficiency; a detailed analysis of inference performance was beyond its primary scope. Despite these limitations, this work contributes a detailed and fair empirical comparison to the field of BP-free deep learning, particularly by highlighting the potential of MF as a practical, high-performance, and energy-efficient algorithm for applications based on MLPs.